
Dvanáctý případ užití je specifikován v tabulce č. \ref{tab:usecase:UC12}.

\begin{UseCaseTable}
  {Filtrování dotazů na čtecí API}
  {UC12}
  {F12}
  {Získání pouze té podmnožiny pozorování a metadat, která odpovídá zadaným filtračním kritériím}
  {Vývojář aplikací třetí strany}
  {}
  {Žádné -- API je veřejně dostupné}
  {Konzument získal odfiltrovanou podmnožinu pozorování a metadat odpovídající zadaným kritériím}

\UseCaseScenario{Základní scénář}
\UseCaseStep{vývojář}{Odešle systému požadavek na čtecí API s parametry pro filtraci (časový interval, geografická oblast, identifikátor subjektu zájmu, senzoru či veličiny).}
\UseCaseStep{systém}{Ověří, že zadané filtrační parametry jsou v očekávaném tvaru a hodnoty z číselníků odpovídají povoleným položkám.}
\UseCaseStep{systém}{Aplikuje filtry na evidenci pozorování a metadat.}
\UseCaseStep{systém}{Vrátí odpověď obsahující pouze záznamy odpovídající všem filtračním kritériím.}

\UseCaseScenario{Alternativní scénář -- neplatný filtrační parametr}
\UseCaseStepCustom{A1}{vývojář}{shodné s krokem 1 základního scénáře.}
\UseCaseStepCustom{A2}{systém}{Zjistí, že některý z filtračních parametrů má neplatný formát nebo odkazuje na neexistující entitu.}
\UseCaseStepCustom{A3}{systém}{Vrátí chybovou odpověď s odpovídajícím HTTP stavovým kódem a popisem problematického parametru.}

\UseCaseScenario{Alternativní scénář -- prázdný výsledek}
\UseCaseStepCustom{B1}{vývojář, systém}{shodné s kroky 1--3 základního scénáře.}
\UseCaseStepCustom{B2}{systém}{Zjistí, že filtračním kritériím neodpovídá žádný záznam.}
\UseCaseStepCustom{B3}{systém}{Vrátí prázdnou množinu výsledků v očekávané struktuře odpovědi.}

\end{UseCaseTable}
